\documentclass[12pt]{article}
\usepackage[utf8]{inputenc}
\usepackage[T1]{fontenc}
\usepackage[margin=1in]{geometry}
\usepackage{setspace}
\usepackage{xcolor}
\usepackage[colorlinks=true, urlcolor=blue!60!black, linkcolor=blue!60!black]{hyperref}
\setlength{\parindent}{2em}
\setlength{\parskip}{1.25em}
\renewcommand{\baselinestretch}{1.0}
\begin{document}
\begin{center}
{\LARGE \textit{Reynolds v. United States}, 98 U.S. 145 (1879)\par}
\vspace{0.8em}
{\large Case Brief\par}
\end{center}
\vspace{1em}

\section{Memory Jogger}

In its first major interpretation of the First Amendment's Free Exercise Clause, the Supreme Court drew a sharp line between religious belief and religious conduct, holding that sincere religious duty does not excuse violating a generally applicable federal criminal prohibition on polygamy (\href{https://www.courtlistener.com/opinion/89858/reynolds-v-united-states/}{CourtListener}; \href{https://supreme.justia.com/cases/federal/us/98/145/}{Justia}).

\section{Detailed Case Facts}

George Reynolds, a resident of the Territory of Utah and a devout member of the Church of Jesus Christ of Latter-day Saints (commonly known as the Mormon Church), was married to Mary Ann Tuddenham (\href{https://supreme.justia.com/cases/federal/us/98/145/}{Justia}). While his first wife was still living, Reynolds entered into a second marriage with Amelia Jane Schofield in a ceremony performed by Daniel H. Wells, an official recognized by the Church as having authority to solemnize marriages, with the permission of Church authorities (\href{https://supreme.justia.com/cases/federal/us/98/145/}{Justia}).

At trial, Reynolds proved that it was an accepted doctrine of his Church that male members had a religious duty, circumstances permitting, to practice polygamy; that this duty was believed to be directly enjoined by Almighty God in a revelation to Joseph Smith; and that refusing to practice polygamy when circumstances permitted would be punished by ``damnation in the life to come'' (\href{https://supreme.justia.com/cases/federal/us/98/145/}{Justia}). However, Congress had enacted the Morrill Anti-Bigamy Act of 1862 (12 Stat. 501), codified as Revised Statutes \S~5352, which provided that every person having a living husband or wife who married another in a Territory or other place under the exclusive jurisdiction of the United States was guilty of bigamy, punishable by a fine of not more than \$500 and imprisonment for not more than five years (\href{https://supreme.justia.com/cases/federal/us/98/145/}{Justia}).

Several evidentiary and jury-selection disputes also arose during Reynolds's prosecution. First, the indictment was returned by a fifteen-person grand jury impaneled under an 1870 Utah territorial statute (Compiled Laws of Utah 1876, p.~357, \S~4), whereas federal Revised Statutes \S~808 required grand juries in United States circuit and district courts to consist of sixteen to twenty-three persons (\href{https://supreme.justia.com/cases/federal/us/98/145/}{Justia}). Second, during \textit{voir dire}, the trial court overruled Reynolds's challenges for cause to prospective petit jurors Eli Ransohoff and Charles Read, both of whom admitted having formed impressions or opinions about the case from outside sources (Read had not expressed his opinion and stated he did not think it would influence his verdict), while sustaining the prosecution's challenges to prospective jurors Homer Brown and John W. Snell after they refused to answer whether they were living in polygamy on self-incrimination grounds (\href{https://supreme.justia.com/cases/federal/us/98/145/}{Justia}).

Third, when the prosecution sought to prove the second marriage, Deputy Marshal Arthur Pratt was unable to serve a subpoena on Amelia Jane Schofield. An initial subpoena mistakenly listed her name as ``Mary Jane Schobold,'' and when Pratt went to Reynolds's home and asked where she was, Reynolds replied that she was not home, told Pratt ``that will be for you to find out,'' and stated words to the effect that she would not appear in the case (\href{https://supreme.justia.com/cases/federal/us/98/145/}{Justia}). After a corrected subpoena issued for Amelia Jane Schofield, Pratt returned to Reynolds's house and was told by Reynolds's first wife that Schofield had not been there for two or three weeks. Over Reynolds's Sixth Amendment Confrontation Clause objection, the trial court permitted A. S. Patterson and other witnesses to testify to what Schofield had sworn at an earlier trial of Reynolds for the same offense under a different indictment (\href{https://supreme.justia.com/cases/federal/us/98/145/}{Justia}).

Finally, the trial judge refused Reynolds's requested jury instruction that he must be acquitted if he married in conformity with what he believed to be a religious duty. Instead, the judge instructed the jury that deliberately marrying a second time while a first wife was living implied criminal intent regardless of religious belief, and urged jurors to consider the consequences to the ``pure-minded women'' and ``innocent children'' whom the judge described as ``the innocent victims of this delusion'' (\href{https://supreme.justia.com/cases/federal/us/98/145/}{Justia}).

\section{Procedural History}

Reynolds had previously been indicted and tried for the same offense under an earlier indictment returned by a twenty-three-person grand jury, which the Supreme Court of the Territory of Utah held invalid under territorial law (\textit{United States v. Reynolds}, 1 Utah 226) (\href{https://supreme.justia.com/cases/federal/us/98/145/}{Justia}). He was subsequently re-indicted by a fifteen-person grand jury in the District Court for the Third Judicial District of the Territory of Utah under Revised Statutes \S~5352 (\href{https://supreme.justia.com/cases/federal/us/98/145/}{Justia}).

Reynolds pleaded in abatement that the fifteen-person grand jury was illegal under Revised Statutes \S~808; the District Court overruled the plea, and Reynolds pleaded not guilty. Following trial, the jury convicted him of bigamy, and the District Court sentenced him to two years' imprisonment at hard labor and a \$500 fine (\href{https://supreme.justia.com/cases/federal/us/98/145/}{Justia}). The Supreme Court of the Territory of Utah affirmed the conviction and sentence. Reynolds then sued out a writ of error to the Supreme Court of the United States, where the case was argued on November 14--15, 1878, and decided on January 6, 1879 (\href{https://www.oyez.org/cases/1850-1900/98us145}{Oyez}; \href{https://supreme.justia.com/cases/federal/us/98/145/}{Justia}).

Although the Supreme Court initially affirmed the judgment in full, Reynolds later filed a petition for rehearing during the same term pointing out that Revised Statutes \S~5352 authorized imprisonment only and did not authorize imprisonment ``at hard labor'' (\href{https://supreme.justia.com/cases/federal/us/98/145/}{Justia}). On rehearing, Chief Justice Waite issued a supplemental opinion vacating the prior affirmance, reversing the judgment solely to correct the form of the sentence, and remanding with instructions to enter a new sentence identical to the first except omitting hard labor (\href{https://supreme.justia.com/cases/federal/us/98/145/}{Justia}).

\section{Judicial Votes}

\begin{itemize}
\item \textbf{Vote Breakdown:} 9--0 on the Free Exercise Clause and five of the six assignments of error; 8--1 on the Confrontation Clause evidentiary issue regarding Amelia Jane Schofield's prior testimony (\href{https://supreme.justia.com/cases/federal/us/98/145/}{Justia}; \href{https://constitutioncenter.org/the-constitution/supreme-court-case-library/reynolds-v-united-states}{National Constitution Center}).
\item \textbf{Majority Opinion (8 Justices in full):} Chief Justice Morrison R. Waite authored the opinion of the Court, joined in full by Justices Nathan Clifford, Noah Haynes Swayne, Samuel Freeman Miller, William Strong, Joseph P. Bradley, Ward Hunt, and John Marshall Harlan (\href{https://www.oyez.org/cases/1850-1900/98us145}{Oyez}; \href{https://supreme.justia.com/cases/federal/us/98/145/}{Justia}).
\item \textbf{Concurring in Part and Dissenting in Part (1 Justice):} Justice Stephen Johnson Field concurred with the majority on all points---including the First Amendment Free Exercise holding---except the admission of Amelia Jane Schofield's former testimony, from which he dissented (\href{https://supreme.justia.com/cases/federal/us/98/145/}{Justia}; \href{https://constitutioncenter.org/the-constitution/supreme-court-case-library/reynolds-v-united-states}{National Constitution Center}).
\end{itemize}

\section{Holding}

\textbf{Judgment:} Initially affirmed (98 U.S. at 168); on rehearing later in the same term, vacated and reversed solely as to the ``at hard labor'' portion of the sentence, and remanded for entry of a corrected sentence of two years' imprisonment and a \$500 fine (\href{https://supreme.justia.com/cases/federal/us/98/145/}{Justia}).

\textbf{Rule Announced:} Under the Free Exercise Clause of the First Amendment, Congress is deprived of all legislative power over mere religious belief and opinion, but remains free to regulate or criminalize overt actions that violate social duties or are subversive of good order; accordingly, a defendant's sincere religious belief or religiously mandated duty cannot serve as a defense to a generally applicable criminal statute prohibiting bigamy (\href{https://supreme.justia.com/cases/federal/us/98/145/}{Justia}). On the accompanying procedural issues, the Court held that: (1) Utah territorial law governed the size of territorial grand juries; (2) a trial court's ruling on a juror challenge for cause based on a non-fixed opinion will not be reversed absent manifest error; (3) excluding practicing polygamists from a bigamy jury was proper; (4) under the forfeiture-by-wrongdoing doctrine, a defendant who wrongfully procures a witness's absence forfeits his Sixth Amendment Confrontation Clause objection to admitting the witness's former testimony; and (5) the trial court's jury charge regarding the social harms of polygamy was not reversible error (\href{https://supreme.justia.com/cases/federal/us/98/145/}{Justia}).

\textbf{Holding vs. Dicta:} The binding Free Exercise holding is that the First Amendment does not entitle religious adherents to an exemption from a valid, generally applicable criminal prohibition on polygamy, and that religious duty does not negate criminal intent when the defendant knowingly commits the prohibited act (\href{https://supreme.justia.com/cases/federal/us/98/145/}{Justia}). By contrast, Chief Justice Waite's sweeping historical generalizations---including his invocation of Thomas Jefferson's 1802 Danbury Baptist letter ``wall of separation'' metaphor as ``almost as an authoritative declaration of the scope and effect of the amendment,'' his remarks that polygamy was ``odious among the northern and western nations of Europe'' and historically ``almost exclusively a feature of the life of Asiatic and of African people,'' his citation of Francis Lieber's political theory that polygamy ``fetters the people in stationary despotism,'' and his hypothetical examples of human sacrifice and widow-burning (\textit{sati})---constitute influential dicta rather than strictly necessary elements of the judgment (\href{https://supreme.justia.com/cases/federal/us/98/145/}{Justia}).

\section{Analysis of Opinions}

\subsection{Majority Opinion (Chief Justice Waite)}

Chief Justice Waite addressed Reynolds's six assignments of error in turn, reserving his most consequential constitutional analysis for the fifth issue: the Free Exercise defense (\href{https://supreme.justia.com/cases/federal/us/98/145/}{Justia}).

On the four threshold procedural challenges, Waite first held that territorial district courts are not Article III ``circuit and district courts of the United States'' under Revised Statutes \S~808, so Utah's fifteen-person grand jury statute controlled. Second, relying on Lord Coke and Chief Justice Marshall's standard from \textit{Burr's Trial}, Waite concluded that juror Charles Read's unexpressed, hypothetical impression did not establish manifest partiality requiring reversal of the trial court's factual determination. Third, Waite upheld the exclusion of prospective jurors who were living in polygamy, reasoning that they could not be impartial and that any technical distinction between challenges for cause and for favor was immaterial. Fourth, invoking English common law dating back to \textit{Lord Morley's Case} (1666) and the equitable maxim that no one may profit from his own wrong, Waite held that the Sixth Amendment Confrontation Clause does not bar introducing a witness's prior testimony from a former trial between the same parties on the same issue when the accused wrongfully procures the witness's absence; because Reynolds's statements to Deputy Marshal Pratt shifted the burden to him to explain Schofield's absence and he declined to do so under oath, the trial court did not manifestly err in admitting her prior testimony (\href{https://supreme.justia.com/cases/federal/us/98/145/}{Justia}).

Turning to the First Amendment, Waite confirmed that the Free Exercise Clause binds Congress when legislating for federal Territories. Because the Constitution does not define ``religion,'' Waite looked to founding-era history in Virginia: James Madison's \textit{Memorial and Remonstrance} (1785), Thomas Jefferson's Virginia Statute for Establishing Religious Freedom (1786)---which declared that civil government may interfere ``when principles break out into overt acts against peace and good order''---and Jefferson's 1802 letter to the Danbury Baptist Association stating that legislative powers ``reach actions only, and not opinions'' and that the First Amendment built ``a wall of separation between church and State'' (\href{https://supreme.justia.com/cases/federal/us/98/145/}{Justia}). Treating Jefferson's formulation as an authoritative statement of the Clause's scope, Waite announced the belief-action dichotomy: Congress cannot regulate religious opinions, but may prohibit actions that violate social duties or subvert good order (\href{https://supreme.justia.com/cases/federal/us/98/145/}{Justia}; \href{https://firstamendment.mtsu.edu/article/reynolds-v-united-states/}{First Amendment Encyclopedia}).

Waite then examined the legal history of polygamy, noting that it had been criminalized in England since the Statute of 1 James I, c.~11 (punishable by death), reenacted across the American colonies, and reenacted by Virginia in December 1788 shortly after adopting Jefferson's religious-freedom statute and recommending a federal bill of rights (\href{https://supreme.justia.com/cases/federal/us/98/145/}{Justia}). Characterizing marriage as both a sacred obligation and a civil contract upon which society is built, Waite quoted political theorist Francis Lieber and Chancellor James Kent for the proposition that polygamy fosters ``the patriarchal principle'' and ``stationary despotism,'' whereas monogamy supports democratic governance. Rejecting Reynolds's claim for a religious exemption, Waite analogized religiously mandated polygamy to religiously mandated human sacrifice or widow-burning on a funeral pyre, warning that excusing criminal conduct on the basis of religious belief would ``make the professed doctrines of religious belief superior to the law of the land, and, in effect, to permit every citizen to become a law unto himself'' (\href{https://supreme.justia.com/cases/federal/us/98/145/}{Justia}). Waite also rejected Reynolds's \textit{mens rea} argument: because Reynolds knowingly married a second time while his first wife lived and knew the law forbade it, criminal intent was implied, and religious belief that the law was invalid could not negate intent. Finally, Waite upheld the trial judge's remarks about the ``innocent victims'' of polygamy as a permissible reminder of the jury's duty rather than an improper appeal to passion or prejudice (\href{https://supreme.justia.com/cases/federal/us/98/145/}{Justia}).

\subsection{Concurrence in Part and Dissent in Part (Justice Field)}

Justice Stephen J. Field wrote a brief, four-sentence separate opinion. He joined the majority on all constitutional and procedural questions---including the First Amendment Free Exercise holding upholding Revised Statutes \S~5352---but dissented solely from the admission of Amelia Jane Schofield's former testimony (\href{https://supreme.justia.com/cases/federal/us/98/145/}{Justia}). Field argued that the prosecution had failed to lay a sufficient evidentiary foundation showing that Reynolds had wrongfully procured Schofield's absence, asserting that the very common-law authorities cited by Chief Justice Waite ``establish conclusively the exact reverse'' (\href{https://supreme.justia.com/cases/federal/us/98/145/}{Justia}).

\section{Examples: Future Applications}

\begin{enumerate}
\item \textbf{Same-Side (Government Prevails):} \textit{Religiously Motivated Use of Controlled Substances.} Members of a religious sect ingest an illegal hallucinogen as a central sacrament and challenge their criminal prosecution (or resulting denial of unemployment benefits) under a neutral, generally applicable state drug statute. Under \textit{Reynolds}'s belief-action dichotomy, the state may not criminalize belief in the sacrament's spiritual value, but it may prohibit the overt act of ingesting a controlled substance without granting a constitutional exemption---the exact reasoning the Court later revived in \href{https://supreme.justia.com/cases/federal/us/494/872/}{\textit{Employment Division v. Smith}, 494 U.S. 872 (1990)}.

\item \textbf{Same-Side (Government Prevails):} \textit{Religious Objection to Mandatory Social Security Taxes.} An employer whose faith teaches that community members must care for their own elderly and forbids participating in government insurance programs refuses to pay federal Social Security payroll taxes. Because the tax law is a neutral, generally applicable regulation of economic conduct essential to social order rather than a regulation of belief, \textit{Reynolds} dictates that religious duty cannot excuse noncompliance---as the Court later held in \href{https://supreme.justia.com/cases/federal/us/455/252/}{\textit{United States v. Lee}, 455 U.S. 252 (1982)}.

\item \textbf{Opposite-Side (Religious Claimant Prevails):} \textit{Criminalizing a Religious Belief or Creed.} A federal territorial statute makes it a misdemeanor to profess belief in plural marriage or to preach that polygamy is divinely ordained, even without entering a second marriage. Under \textit{Reynolds}, Congress has zero legislative authority over ``mere religious belief and opinions'' (\href{https://supreme.justia.com/cases/federal/us/98/145/}{Justia}); punishing pure belief or theological advocacy crosses directly into the forbidden sphere of opinion rather than overt acts against public order.

\item \textbf{Opposite-Side (Religious Claimant Prevails):} \textit{Targeted Prohibition of a Minority Religious Ritual.} A city council passes ordinances specifically drafted to ban ritual animal sacrifice practiced by a single minority faith, while exempting secular slaughter, hunting, and kosher butchery. Even under \textit{Reynolds} and its modern successor \href{https://supreme.justia.com/cases/federal/us/494/872/}{\textit{Smith}}, the government may not target a religious practice because of its religious motivation rather than enforcing a neutral rule of general applicability---the principle enforced in \href{https://supreme.justia.com/cases/federal/us/508/520/}{\textit{Church of the Lukumi Babalu Aye, Inc. v. City of Hialeah}, 508 U.S. 520 (1993)}.

\item \textbf{Fence-Sitter (Close Call):} \textit{Compulsory Secondary Education vs. Insular Religious Community.} A state enforces a neutral, generally applicable compulsory school-attendance law requiring all children to attend school until age sixteen against parents from a centuries-old, self-sufficient agrarian religious community who withdraw their children after eighth grade to preserve their faith and way of life. Under a strict reading of \textit{Reynolds}, refusing to send children to school is an overt act violating a valid, neutral law of social organization, so the state would prevail. Yet because the conduct involves no violent harm, threatens no broad breakdown of social order, and implicates parental upbringing alongside deeply rooted religious survival, the Supreme Court departed from rigid \textit{Reynolds} formalism in \href{https://supreme.justia.com/cases/federal/us/406/205/}{\textit{Wisconsin v. Yoder}, 406 U.S. 205 (1972)}, demonstrating the enduring tension between \textit{Reynolds} and accommodationist scrutiny.
\end{enumerate}

\section{Critique}

\textit{Reynolds v. United States} remains one of the most fiercely debated decisions in American constitutional law, sitting at the center of a century-and-a-half struggle over whether the Free Exercise Clause requires religious exemptions from neutral laws.

\textbf{Originalist and Textualist Critiques:} Leading originalist scholar Michael W. McConnell has sharply criticized Chief Justice Waite's historical methodology (\href{https://law.yale.edu/sites/default/files/area/clinic/document/michael_mcconnell_the_origins_and_historical_understanding_of_free_exercise_of_religion_103_harv._l._rev._1409_1989.pdf}{McConnell, 103 Harv. L. Rev. 1409 (1990)}). McConnell argues that Waite erred by treating Thomas Jefferson---who was in France during the Constitutional Convention and the drafting of the Bill of Rights---as the authoritative voice of the First Amendment, while overlooking James Madison's broader conception of religious duty and early state constitutional provisos that protected religious exercise unless it disturbed the peace or safety of the state. Textualists also note that the First Amendment protects the ``free \textit{exercise}'' of religion, a word denoting active practice rather than purely internal belief. More recent scholarship, such as Clark B. Lombardi's examination of \textit{Reynolds} after \href{https://www.law.cornell.edu/supremecourt/text/19-123}{\textit{Fulton v. City of Philadelphia}, 593 U.S. 522 (2021)}, emphasizes that even Waite's own opinion did not claim all religious conduct was unprotected, but rather asked whether polygamy fell within the traditional limits of peace and good order (\href{https://law.ua.edu/wp-content/uploads/2024/06/Lombardi.pdf}{Lombardi, 75 Ala. L. Rev. 1009 (2024)}).

\textbf{Progressive, Doctrinal, and Historical Critiques:} Doctrinal and historical critics observe that \textit{Reynolds} reflects nineteenth-century Protestant majoritarianism and anti-Mormon hostility rather than neutral constitutional principle. Waite's characterization of polygamy as ``odious'' to European civilization and associated with ``Asiatic and of African people'' and ``stationary despotism'' relied on culturally biased assumptions rather than empirical proof of harm (\href{https://supreme.justia.com/cases/federal/us/98/145/}{Justia}; \href{https://firstamendment.mtsu.edu/article/reynolds-v-united-states/}{First Amendment Encyclopedia}). Doctrinally, the Court temporarily moved away from \textit{Reynolds}'s categorical rejection of exemptions during the Warren and Burger Courts, adopting compelling-interest strict scrutiny in \href{https://supreme.justia.com/cases/federal/us/374/398/}{\textit{Sherbert v. Verner}, 374 U.S. 398 (1963)}, and \href{https://supreme.justia.com/cases/federal/us/406/205/}{\textit{Wisconsin v. Yoder}, 406 U.S. 205 (1972)}, before Justice Scalia revived \textit{Reynolds}'s ``law unto himself'' warning in \href{https://supreme.justia.com/cases/federal/us/494/872/}{\textit{Employment Division v. Smith}, 494 U.S. 872 (1990)} (\href{https://harvardlawreview.org/print/vol-137/a-law-unto-himselfemp-div-v-smith-494-u-s-872-879-1990-quoting-reynolds-v-united-states-98-u-s-145-167-1879-free-exercise-unequal-value-and-the-future-of-public/}{137 Harv. L. Rev. (2024)}). Conversely, defenders of \textit{Reynolds} and \textit{Smith}---including many progressive and egalitarian scholars---argue that granting constitutional exemptions from generally applicable laws risks undermining civil rights, public health, and the equal rule of law by allowing religious objections to override democratically enacted protections for third parties.

\textbf{Underlying Competing Values:} At bottom, \textit{Reynolds} exposes the irreducible tension between \textit{freedom of religious conscience} (which for most faiths necessarily entails acting on divine commands, not merely holding private thoughts) and \textit{democratic rule of law and equal citizenship} (which resists allowing individuals to opt out of neutral laws that protect social order and third-party rights).

\section{Key Quotations}

\begin{quote}
``Congress was deprived of all legislative power over mere opinion, but was left free to reach actions which were in violation of social duties or subversive of good order.''
\end{quote}
\noindent --- Chief Justice Morrison R. Waite, \textit{Reynolds v. United States}, 98 U.S. 145, 164 (1879) (\href{https://supreme.justia.com/cases/federal/us/98/145/}{Justia}; \href{https://www.courtlistener.com/opinion/89858/reynolds-v-united-states/}{CourtListener}).

\begin{quote}
``Laws are made for the government of actions, and while they cannot interfere with mere religious belief and opinions, they may with practices.''
\end{quote}
\noindent --- Chief Justice Morrison R. Waite, \textit{Reynolds v. United States}, 98 U.S. 145, 166 (1879) (\href{https://supreme.justia.com/cases/federal/us/98/145/}{Justia}; \href{https://www.courtlistener.com/opinion/89858/reynolds-v-united-states/}{CourtListener}).

\begin{quote}
``Can a man excuse his practices to the contrary because of his religious belief? To permit this would be to make the professed doctrines of religious belief superior to the law of the land, and, in effect, to permit every citizen to become a law unto himself. Government could exist only in name under such circumstances.''
\end{quote}
\noindent --- Chief Justice Morrison R. Waite, \textit{Reynolds v. United States}, 98 U.S. 145, 166--167 (1879) (\href{https://supreme.justia.com/cases/federal/us/98/145/}{Justia}; \href{https://www.courtlistener.com/opinion/89858/reynolds-v-united-states/}{CourtListener}).

\begin{quote}
``Marriage, while from its very nature a sacred obligation, is nevertheless, in most civilized nations, a civil contract, and usually regulated by law.''
\end{quote}
\noindent --- Chief Justice Morrison R. Waite, \textit{Reynolds v. United States}, 98 U.S. 145, 165 (1879) (\href{https://supreme.justia.com/cases/federal/us/98/145/}{Justia}; \href{https://www.courtlistener.com/opinion/89858/reynolds-v-united-states/}{CourtListener}).

\begin{quote}
``I concur with the majority of the court on the several points decided except one --- that which relates to the admission of the testimony of Amelia Jane Schofield given on a former trial upon a different indictment. I do not think that a sufficient foundation was laid for its introduction. The authorities cited by the Chief Justice to sustain its admissibility seem to me to establish conclusively the exact reverse.''
\end{quote}
\noindent --- Justice Stephen J. Field, concurring in part and dissenting in part, \textit{Reynolds v. United States}, 98 U.S. 145, 168 (1879) (\href{https://supreme.justia.com/cases/federal/us/98/145/}{Justia}; \href{https://www.courtlistener.com/opinion/89858/reynolds-v-united-states/}{CourtListener}).

\end{document}
